\documentclass[11pt]{article}

\usepackage[margin=0.72in]{geometry}
\usepackage{amsmath,amssymb,amsthm}
\usepackage{hyperref}
\usepackage{booktabs}

\hypersetup{colorlinks=true,linkcolor=blue,citecolor=blue,urlcolor=blue}

\newtheorem{theorem}{Theorem}
\newtheorem{proposition}[theorem]{Proposition}
\newtheorem{corollary}[theorem]{Corollary}

\newcommand{\Tr}{\operatorname{Tr}}
\newcommand{\Sym}{\operatorname{Sym}}

\title{The Torsion Sector of the Tetrad Connection\\
as the Regular Representation of $S_4$}
\author{Stephen Nelson\\
\small Field Dynamics Research\\
\small\texttt{stephen@fielddynamics.org}}
\date{}

\begin{document}
\maketitle

\begin{abstract}
The affine connection on the standard tetrad frame has $K^3=64$
independent components, of which 40 are symmetric and 24 are
antisymmetric in the lower index pair.  General relativity's torsion-free condition removes the
associated 24-dimensional torsion sector.
This note shows that, at the physical frame dimension $K=4$,
the torsion sector $V\otimes\Lambda^2V^*$, equivalently
$V\otimes\Lambda^2V$ under the self-duality of the permutation
representation, carries the regular
representation of the symmetric group $S_4$ acting by
permutation of the four frame legs.  The regular representation
is the richest possible finite-group module: it contains every
irreducible $S_4$-representation at full multiplicity.  Among the
five $S_4$ irreps, the sign representation (the unique scalar
parity detector) appears with multiplicity one in the torsion
sector and multiplicity zero in the Christoffel sector.  General
relativity's torsion-free condition therefore removes the only
scalar mechanism for distinguishing even from odd frame
permutations.  The identification holds
at $K=4$ and at the degenerate $K=2$; at every other frame
dimension the torsion count $K^2(K-1)/2$ does not equal $K!$.
\end{abstract}


\section{The torsion question}

The vielbein formulation of general relativity describes the
gravitational field through a frame field $e^a{}_\mu$ with $K=4$
internal frame legs~\cite{Wald1984,Hehl1976}.  The affine
connection $\Gamma^a{}_{bc}$ on this frame carries three frame
indices and therefore has $K^3=64$ independent components.  Under
exchange of the lower index pair, these components divide into two
families: 40 in the lower-index symmetric sector, corresponding to
the torsion-free connection, and 24 in the lower-index
antisymmetric sector, sampled by the torsion
tensor~\cite{Hehl1976}.

Einstein formulated general relativity in 1915 with the assumption
that the connection is symmetric in its lower indices, setting all
24 torsion components to zero.  Cartan~\cite{Cartan1922} showed in
1922 that the mathematical framework accommodates torsion naturally,
and Hehl and collaborators~\cite{Hehl1976} developed the full
Einstein-Cartan theory, in which the 24 torsion components couple
to the spin density of matter.  The resulting theory reproduces
every prediction of general relativity in vacuum and differs only
inside extremely dense spinning matter.

The torsion sector has been studied extensively under the continuous
Lorentz group.  Its decomposition into irreducible Lorentz
representations (vector, axial vector, and tensor parts, with
$4+4+16=24$ components) is standard~\cite{Hehl1976}.  What has
received less attention is the discrete symmetry structure of the
torsion sector under the permutation group of the frame legs
themselves.  The four legs of the tetrad can be permuted in $4!=24$
ways, and the symmetric group $S_4$ acts naturally on any tensor
built from the frame space $V$.  The question addressed here is
whether this action reveals additional algebraic structure in the
torsion sector, beyond what the continuous Lorentz decomposition
provides.

The answer, as I show below, is that it does.  The $S_4$ action on
the torsion sector $V\otimes\Lambda^2V^*$ (equivalently
$V\otimes\Lambda^2V$ as an $S_4$-module) has a particularly clean
character: it is the regular representation of $S_4$.  This is the
strongest possible finite-group structure a $|G|$-dimensional space
can carry; it contains every irreducible representation at full
multiplicity.  The result holds at $K=4$ and at the
degenerate case $K=2$; it fails at every other positive integer,
because the numerical coincidence $K^2(K-1)/2=K!$ is satisfied
only at these two values.  The physical spacetime dimension is
therefore distinguished, within the tetrad formalism, as the unique
nontrivial dimension where the torsion sector exhausts the full
permutation structure of the frame.

In what follows, Section~2 recalls the $S_4$ action on the frame
space and its exterior powers, Section~3 computes the character of
the induced representation on $V\otimes\Lambda^2V^*$, Section~4
establishes the regular-representation identification, and
Section~5 shows that the sign representation lives exclusively in
the torsion sector.  The notation is standard: $V$ is a real $K$-dimensional
vector space, $S_4$ acts by permutation of a chosen basis, and
characters are traces of representation matrices.


\section{The frame permutation group}

The connection $\Gamma^a{}_{bc}$ carries three indices from the set
$\{0,1,2,3\}$.  Any tensor whose indices run over a set of size
$K$ inherits a canonical $S_K$ action by permutation of the index
labels.  This action exists prior to and independently of any inner
product on $V$; the Lorentz group $O(1,3)$ arises only after the
inner product $\eta$ is imposed, restricting $GL(4)$ to its
$\eta$-preserving subgroup.

The $S_4$ representation theory computed here describes the
index-permutation module of the torsion component space.  Adding
$\eta$ breaks $S_4$ to the subgroup that preserves the signature
pattern (essentially $S_1\times S_3$, permuting the three
spacelike legs while fixing the timelike leg), but the
representation-theoretic content of the full $S_4$ action remains
as the pre-metric algebraic structure of the frame.  The Lorentz
decomposition of torsion into vector, axial vector, and tensor
parts ($4+4+16$) is a post-$\eta$ decomposition of the same
24-dimensional component space.  The regular representation found
below is a pre-$\eta$ index-permutation structure on that space.
The relationship between these two decompositions is left open.

The character $\chi_V(\sigma)$ used throughout this note counts
fixed frame labels; it does not evaluate any inner product.  The
computation is therefore independent of whether the frame carries
Lorentzian, Euclidean, or no signature at all.

Let $\{e_0,e_1,e_2,e_3\}$ be a basis for the four-dimensional
frame space $V$.  The symmetric group $S_4$ acts on $V$ by
permuting these basis vectors: for $\sigma\in S_4$,
\[
        \sigma\cdot e_a = e_{\sigma(a)}.
\]
This is the defining (or natural) representation of $S_4$ on $V$.
Its character is the number of fixed points:
\[
        \chi_V(\sigma) = |\{a : \sigma(a)=a\}|.
\]

The action extends to tensor products and exterior powers by
linearity.  On the exterior square $\Lambda^2V$, whose basis
elements are the antisymmetric wedge products $e_i\wedge e_j$
$(i<j)$, the group acts by
\[
        \sigma\cdot(e_i\wedge e_j)
        = e_{\sigma(i)}\wedge e_{\sigma(j)}.
\]
The character of this representation is given by the standard
formula~\cite{FultonHarris}
\[
        \chi_{\Lambda^2V}(\sigma)
        = \frac{\chi_V(\sigma)^2 - \chi_V(\sigma^2)}{2}.
\]

The torsion sector of the connection is the tensor product
$V\otimes\Lambda^2V^*$.  For the permutation representation of
$S_4$, the natural representation is self-dual ($V^*\cong V$ as
$S_4$-modules), so the character computation may be carried out
equivalently on $V\otimes\Lambda^2V$ without changing the result.
The dimension is $K\times K(K-1)/2 = 4\times 6 = 24$.  Each
torsion component $T^a{}_{[bc]}$ is indexed by one frame label $a$
from the set $\{0,1,2,3\}$ and one antisymmetric pair $[bc]$ from
the $\binom{4}{2}=6$ unordered pairs.  This is the incidence
structure of a combinatorial tetrahedron: four vertices connected
by six edges, giving $4\times 6=24$ vertex-edge incidences.  The
same count gives $|S_4|=4!=24$, which is why the regular
representation fits exactly into this component space.

The $S_4$ action on this space is the tensor
product of the two actions above, with character
\[
        \chi_{V\otimes\Lambda^2V}(\sigma)
        = \chi_V(\sigma)\cdot\chi_{\Lambda^2V}(\sigma).
\]


\section{The character computation}

The group $S_4$ has five conjugacy classes, distinguished by cycle
type.  For each class, one computes $\chi_V$, then $\chi_V$ of the
squared element $\sigma^2$, then $\chi_{\Lambda^2V}$, and finally
the product.

\begin{center}
\begin{tabular}{lccccc}
\toprule
Cycle type & Size & $\chi_V(\sigma)$ & $\chi_V(\sigma^2)$
           & $\chi_{\Lambda^2V}(\sigma)$
           & $\chi_{V\otimes\Lambda^2V}(\sigma)$ \\
\midrule
$(1^4)$: identity         & 1  & 4 & 4 & 6  & 24 \\
$(2,1^2)$: transpositions & 6  & 2 & 4 & 0  & 0  \\
$(3,1)$: 3-cycles         & 8  & 1 & 1 & 0  & 0  \\
$(4)$: 4-cycles           & 6  & 0 & 0 & 0  & 0  \\
$(2^2)$: double transpos. & 3  & 0 & 4 & $-2$ & 0  \\
\bottomrule
\end{tabular}
\end{center}

The entries are computed as follows.

\textbf{Identity.}  All four legs are fixed, so $\chi_V(e)=4$.
Then $\chi_V(e^2)=4$ and $\chi_{\Lambda^2V}(e)=(16-4)/2=6$.
Product: $4\times 6=24$.

\textbf{Transpositions.}  A transposition fixes two legs,
so $\chi_V=2$.  Its square is the identity, so $\chi_V(\sigma^2)=4$.
Then $\chi_{\Lambda^2V}=(4-4)/2=0$.  Product: $2\times 0=0$.

\textbf{3-cycles.}  A 3-cycle fixes one leg, so $\chi_V=1$.
Its square is another 3-cycle (same conjugacy class), so
$\chi_V(\sigma^2)=1$.  Then $\chi_{\Lambda^2V}=(1-1)/2=0$.
Product: $1\times 0=0$.

\textbf{4-cycles.}  A 4-cycle fixes no legs, so $\chi_V=0$.
Its square is a double transposition, so $\chi_V(\sigma^2)=0$.
Then $\chi_{\Lambda^2V}=(0-0)/2=0$.  Product: $0\times 0=0$.

\textbf{Double transpositions.}  No legs are fixed, so $\chi_V=0$.
The square of a double transposition is the identity, so
$\chi_V(\sigma^2)=4$.  Then $\chi_{\Lambda^2V}=(0-4)/2=-2$.
Product: $0\times(-2)=0$.


\section{The regular representation}

\begin{theorem}
Let $V$ be the four-dimensional permutation representation of $S_4$.
Then the $S_4$-module $V\otimes\Lambda^2V^*$, equivalently
$V\otimes\Lambda^2V$ under the natural self-duality of the
permutation representation, is isomorphic to the regular
representation of $S_4$.
\end{theorem}

\begin{proof}
A representation whose character equals $|G|$ on the identity and
vanishes on every other group element is, by definition, the
regular representation~\cite{FultonHarris,Serre}.  The table in
Section~3 shows $\chi(e)=24=|S_4|$ and $\chi(\sigma)=0$ for all
$\sigma\neq e$.
\end{proof}

The regular representation is the richest finite-group module
of its dimension.  It decomposes into irreducible representations
as
\[
        \text{reg} = \bigoplus_\rho \dim(\rho)\cdot\rho,
\]
where the sum runs over all irreducible representations $\rho$ of
$S_4$, each appearing with multiplicity equal to its own
dimension.  The irreducible representations of $S_4$ have
dimensions $\{1,1,2,3,3\}$, and the decomposition is confirmed by
\[
        1^2 + 1^2 + 2^2 + 3^2 + 3^2 = 1+1+4+9+9 = 24.
\]
General relativity's torsion-free condition therefore does not
merely remove 24 components from the connection.  It removes a
complete copy of the $S_4$ group algebra, containing every type of
frame-permutation symmetry at full multiplicity.


\section{The sign representation and root-blindness}

Among the five irreducible representations of $S_4$, two are
one-dimensional: the trivial representation, which assigns $+1$ to
every permutation, and the sign representation, which assigns $+1$
to even permutations and $-1$ to odd.  The sign representation is
the unique one-dimensional module that detects frame-orientation
reversal: it is the scalar parity channel of the permutation group.

The regular representation contains every irrep at full
multiplicity, so the torsion sector carries one copy of the sign
representation.  The Christoffel sector, by contrast, carries zero
copies.  This can be verified by computing the inner product of the
Christoffel character $(40,8,1,0,0)$ with the sign character
$(1,-1,1,-1,1)$:
\[
        \langle\chi_{\rm Chris},\chi_{\rm sign}\rangle
        =\frac{1}{24}\bigl[
        1{\cdot}40{\cdot}1
        +6{\cdot}8{\cdot}(-1)
        +8{\cdot}1{\cdot}1
        +6{\cdot}0{\cdot}(-1)
        +3{\cdot}0{\cdot}1
        \bigr]
        =\frac{40-48+8}{24}=0.
\]
The sign representation is entirely absent from the Christoffel
sector.  It lives exclusively in the torsion.

This result gives a representation-theoretic formulation of
root-blindness.  The torsion-free connection, built from the
40-dimensional Christoffel sector alone, has no scalar mechanism
for distinguishing even from odd frame permutations.  The single
copy of the sign representation that would provide such a mechanism
sits in the 24-dimensional torsion sector, which general
relativity's torsion-free condition removes.  In this precise $S_4$-module sense, root-blindness is
not merely an interpretive gloss on the quadratic structure of
the metric; it is the absence of the sign representation from the
Christoffel module.


\section{The uniqueness of $K=4$}

\begin{proposition}
The numerical identity $K^2(K-1)/2 = K!$ holds for positive
integers only at $K\in\{2,4\}$.
\end{proposition}

\begin{proof}
Dividing both sides by $K>0$ gives $K(K-1)/2 = (K-1)!$.  For
$K\geq 3$, the left side grows as $O(K^2)$ while the right side
grows as $O((K-1)!)$, which eventually dominates.  Direct
computation:
\begin{center}
\begin{tabular}{ccc}
\toprule
$K$ & $K^2(K-1)/2$ & $K!$ \\
\midrule
1 & 0 & 1 \\
2 & 2 & 2 \\
3 & 9 & 6 \\
4 & 24 & 24 \\
5 & 50 & 120 \\
6 & 90 & 720 \\
\bottomrule
\end{tabular}
\end{center}
For $K=1$, the left side vanishes while $K!=1$, so no solution
occurs.  For $K\geq 5$, $K! > K^2(K-1)/2$ strictly, since
$(K-1)! > K(K-1)/2$ for $K\geq 5$.  The case $K=2$ is the trivial
two-dimensional frame.  The case $K=4$ is the standard tetrad.
\end{proof}

The regular-representation identification is therefore not a
generic property of frame spaces in arbitrary dimension.  It is
specific to the physical spacetime dimension $K=4$.  At $K=3$, the
torsion sector has 9 components but $|S_3|=6$; the regular
representation does not fit.  At $K=5$, the torsion has 50
components but $|S_5|=120$; again, no match.  The result presented
here is an arithmetical fact about four-dimensional tetrad frames
with no nondegenerate higher-dimensional analogue.


\section{Discussion}

The main observation of this note is that the torsion sector of the
four-dimensional tetrad connection carries algebraic structure that
has not been visible in the standard Lorentz decomposition.  The
$S_4$ regular representation is the strongest possible finite-group
module of dimension 24: it contains every irreducible
representation at full multiplicity.

The Einstein-Cartan programme has traditionally motivated torsion
through its coupling to the spin density of
matter~\cite{Hehl1976}.  The present result offers a complementary
perspective: the torsion sector is not merely a collection of
antisymmetric connection components, but the complete permutation
algebra of the four frame legs.  Setting torsion to zero is
therefore equivalent, in the $S_4$ language, to discarding the
entire regular representation of the frame permutation group.

The irrep decomposition reveals an asymmetry between the two
connection sectors that is not apparent from dimension counting
alone.  The Christoffel sector $V\otimes\Sym^2V$ has character
$\chi_{\rm Chris}=\chi_V\cdot\chi_{\Sym^2V}=(40,8,1,0,0)$ and
decomposes under $S_4$ as
\[
        4\cdot\mathbf{1}
        \oplus 3\cdot V_2
        \oplus 7\cdot V_3
        \oplus 3\cdot (V_3\otimes\mathrm{sgn}),
\]
with zero copies of the sign representation.  Here $V_2$ denotes
the two-dimensional irreducible representation and $V_3$ the
three-dimensional standard representation.  The torsion sector
(24 components), being the regular representation, contains one
copy of every irrep, including the sign.  The sign representation
is the unique one-dimensional $S_4$-module that distinguishes even
from odd frame permutations.  Its complete absence from the
Christoffel sector means that the torsion-free connection lacks the
scalar parity channel of the frame permutation group.

In the companion note~\cite{PhaseA}, the reciprocal partition law
$u(x)=x/(1+x)$ was derived from the $\mathbb Z_2$ exchange
symmetry of the bilinear sector $V\otimes V=\Sym^2V\oplus
\Lambda^2V$.  The present result identifies the torsion sector as
the unique carrier of the sign representation within the full
64-component connection.  Without the torsion, the connection has
no scalar sign channel for detecting even/odd frame-permutation
parity.  With torsion included, the sign representation supplies
exactly this missing scalar parity channel.
The occupation law $u(x)$ from the companion note acts on the
$\Sym^2V/\Lambda^2V$ partition; the present note shows that the
scalar frame-parity channel associated with the antisymmetric side
resides exclusively in the torsion sector that general relativity
removes.

Several directions remain open.  The regular representation
decomposes into the five irreducible $S_4$ types of dimensions
$1,1,2,3,3$, appearing with multiplicities $1,1,2,3,3$
respectively.  How this finite-group decomposition interfaces with
the standard Lorentz decomposition $4+4+16$ remains to be worked
out.  The explicit isomorphism between $V\otimes\Lambda^2V^*$ and
$\mathbb R[S_4]$, and the tensor-component realisation of the
$S_4$ irrep structure, also remain to be constructed.  These
questions are left for future work.


\begin{thebibliography}{9}

\bibitem{Wald1984}
R.~M. Wald,
\emph{General Relativity},
University of Chicago Press (1984).

\bibitem{Hehl1976}
F.~W. Hehl, P.~von der Heyde, G.~D. Kerlick, and J.~M. Nester,
``General relativity with spin and torsion: Foundations and
prospects'',
\emph{Rev.\ Mod.\ Phys.} \textbf{48}, 393--416 (1976).

\bibitem{Cartan1922}
\'E.~Cartan,
``Sur une g\'en\'eralisation de la notion de courbure de Riemann
et les espaces \`a torsion'',
\emph{C.\ R.\ Acad.\ Sci.\ Paris} \textbf{174}, 593--595 (1922).

\bibitem{FultonHarris}
W.~Fulton and J.~Harris,
\emph{Representation Theory: A First Course},
Springer (1991).

\bibitem{Serre}
J.-P.~Serre,
\emph{Linear Representations of Finite Groups},
Springer (1977).

\bibitem{PhaseA}
S.~Nelson,
``An algebraic origin for the simple AQUAL interpolation
function'',
submitted (2026).

\end{thebibliography}

\end{document}
